% discussion.tex
% Discussion: language-audit pass, 2026-07-14 (see
% project_language_audit_2026-07-14 memory). Cut from ~2,700 words to
% ~1,000-1,100: replaced the full per-generator number replay with a
% short synthesis per the audit's suggested text, compressed the
% three per-asset caveats and the multi-asset deferred-questions
% section, moved detail (rejected alternatives, percentile placement)
% to Online Appendix S3, and removed "an honest percentile" and other
% defensive language. No subsections (approved exception was Results
% only); four clearly-separated paragraphs: single-asset result +
% caveats, multi-asset result + caveats, diversification-return bias
% + correction, open extensions. "composer" -> "composition method"
% throughout (conclusion.tex still needs the same swap).

In this study, we found that HMM-WJ improved volatility persistence
relative to HMM-NJ but slightly reduced marginal fit, with pass
rates $2$--$8$ percentage points lower in-sample. Among the alternatives, each
excelled on one property and failed on another for a reason tied to
its own structure: GARCH's conditional-variance dynamics reproduced
temporal clustering but not the marginal's body; the neural baseline
matched some of the temporal structure despite carrying tens of
thousands of parameters, yet still missed the marginal; the
semi-Markov baseline's coarse partition left it unable to span the
full growth-rate range; and bootstrap resampling passed every
distributional test by construction, since it resamples only
observed returns rather than extrapolating beyond the historical
record. Among these alternatives, HMM-WJ was the only non-GARCH,
non-neural model to reduce
the autocorrelation of absolute returns below the shared i.i.d.\
level while retaining most of HMM-NJ's distributional fit, an
ordering that held out-of-sample on the $2025$ evaluation. The model
therefore offers a tradeoff rather than a uniformly better result.

The pass-rate comparisons come with a statistical caveat: two-sample KS
and AD $p$-values assume independent observations, an assumption
the HMM, HSMM, GARCH, and bootstrap generators all violate, so we
weigh pass rates alongside the distribution-level distances ($W_1$,
Hellinger, and the Hill tail index) rather than treat them as
calibrated on a common footing across generators. Three further
caveats apply to the per-asset construction. The $N = 100$ state
resolution, selected by a sensitivity sweep over
$\{30, 60, 90, 100, 150, 200\}$, worked well for the $2{,}766$-day
SPY sample, but shorter histories will support fewer states before
empirical visits per state thin out. The jump parameters
$(\epsilon, \lambda)$ were tuned once against SPY and applied across
the universe; per-ticker tuning on the three cross-asset tickers
produced strong in-sample fits (KS $\ge 98.9\%$, Online
Appendix~\ref{sec:supp-cross-asset}), but is not guaranteed to
transfer to assets with different jump behavior without its own
calibration step. The Student-$t$ emission with $\nu = 5$ improved
kurtosis fidelity over Normal emissions but does not
model skewness directly: the negative skewness both HMM variants
reproduce comes from the asymmetric state-occupancy of the Laplace
partition, so assets whose skew direction the partition cannot
induce will not be matched without adjusting the quantile rule.

We then evaluated the multi-asset extension. The variance
correction reduced the composed marginal's variance error and
improved KS pass rates relative to naive composition, which
inflates variance by $\beta_i^2 \sigma_m^2$, and relative to
Gaussian SIM, which recovers $\beta$ exactly but strips the
marginal entirely. The three
residual-fit methods still had better marginal fit than hybrid
($67$ to $76\%$ KS pass rate versus $50\%$), because fitting a
generator directly to the OLS residual series avoids the
variance-inflation problem rather than correcting for it after the
fact; the choice between the two approaches is governed by whether
the regime-switching and jump structure of the full-return marginal
is itself a property the per-asset generator should have, for
instance when the same generator is used for both single-asset and
multi-asset purposes. On tail risk this tradeoff disappears: hybrid
and the three residual-fit methods all matched on $99\%$ VaR
coverage, differing by less than the cross-ticker spread, so hybrid delivers tail accuracy on par with
a dedicated residual generator without a separate fitting stage,
even though its marginal body fit trails the residual-fit
alternatives.

Several evaluation choices bound these results. We fixed the market
path at the real SPY history, paired Naive and Hybrid on the same
per-ticker JumpHMM innovation path so the comparison isolates the
composition rule (Gaussian SIM and the three residual-fit methods
each drew from their own generators), and skipped the optional
cross-sectional rank reorder; we also evaluated VaR per
ticker rather than at the portfolio level, which the optional
reorder would be needed to support. The variance correction assumes
$\mathrm{Cov}(\tg_i, g_m) \approx 0$, which holds by construction
here since the JumpHMM marginals are fit and sampled independently
of the market path, but would need adjustment for a generator
already correlated with the market (for example, a bootstrap from
joint history). The construction also injects contemporaneous
market dependence only: lagged effects such as lead-lag correlation
or market-driven volatility clustering are not reproduced. Clipping
did not trigger anywhere in this universe because the JumpHMM
marginals had enough variance to keep every ticker's market loading
below the idiosyncratic floor; it becomes relevant for
scenario-driven markets more volatile than the calibration history,
or for universes mixing low-volatility stocks with extreme-beta
assets such as leveraged ETFs. A related, milder attenuation
appeared even without clipping: the hybrid KS pass rate fell from
$67\%$ at low $\beta$ to $44\%$ at high $\beta$, because as the market-driven
component approaches the idiosyncratic floor it accounts for an
increasing share of the marginal, and that component is only as
heavy-tailed as the market path itself.

Independent residuals produce a diversification-return bias under
daily rebalancing: matching each asset's marginal variance does not
restore the contemporaneous cross-asset residual covariance that
common sector and style factors induce in real
markets~\cite{boothFama1992}, so independent residuals overstate the
daily cross-sectional dispersion that rebalancing compounds into a
spurious return premium. This affects every i.i.d.-residual
composition method evaluated here, not only naive composition; block
bootstrap and GARCH(1,1)-$t$ restore each ticker's own temporal
structure but not the cross-asset covariance, so the bias remains
even for the methods that best match each ticker's marginal. In a
$20$-asset target-weight allocator example, the synthetic median
annualized growth rate exceeded the real $2025$ return by about
$22$ percentage points, with a static-versus-daily-rebalance
decomposition consistent with a diversification-return explanation
rather than a calibration
error~\cite{boothFama1992, boucheyNemtchinovWong2015,
markowitz1976longRun} (Online Appendix~\ref{sec:supp-bias}). We
applied a location-shift correction that subtracts a constant drift
in log-wealth from every path, calibrated to match an external
long-run prior~\cite{damodaranERP, ibbotsonSBBI} rather than the
realized test-year outcome. The correction preserves the rank order
of paths and the volatility envelope, but because the drift is
time-dependent in wealth space it does move drawdown dates and
depths, so it recalibrates the ensemble level rather than restoring
the missing cross-asset covariance. Full methodology, the rejected
alternatives, and the percentile placement of the realized year are
given in Online Appendix~\ref{sec:supp-bias}.

Four extensions remain open. Per-ticker tuning of the jump
parameters $(\epsilon, \lambda)$ on each asset's own history would
likely raise marginal pass rates at the cost of a calibration step
per asset. Multi-factor and lagged-factor single-index variants
would replace the single market factor with $K$ factors or lags,
capturing the lagged market dependence the current construction
omits; the variance correction generalizes to both without changing
its algebra, though the empirical evaluation is left to follow-up
work. A jointly specified residual model, such as a synchronized
multivariate block bootstrap, a residual factor model, or a
copula-based dependence model, composed with the single-index factor
instead of independent per-asset draws, would restore cross-asset
residual covariance and address the diversification-return bias
directly; an ordinary per-ticker autoregressive or block-bootstrap
residual model would still be fit and sampled independently for each
ticker, so it would improve each asset's own temporal dependence
without creating dependence between assets. A stress-test harness that
pairs the clipping rule with a market scenario more volatile than
the calibration history would test whether the idiosyncratic floor
keeps per-asset marginals heavy-tailed under exactly the conditions
where naive SIM composition breaks down.
